\begin{papercard}{rethlefsen2021}
\cf{Stages} \stagetag{A} (how apps and sources are searched, filtered, limited, deduplicated and counted), \stagetag{D} (documenting the capture of search results), \stagetag{R} (what to report and where). PRISMA-S (\emph{Systematic Reviews} 10:39, 2021; 19 PDF pages in the appended copy, 16 of text): a 16-item extension to PRISMA for \emph{reporting literature searches}, with an explanation and elaboration for every item. The guidance is written for bibliographic searches, but item wording is generic enough to describe a search of app stores, web sources or grey literature; the card flags what transfers.

\cf{Problem addressed} The search underlies everything downstream: it ``establishes the data available for analysis'' \pc{rethlefsen2021}{1}{the underlying process that estab- lishes the data available for analysis}. Reporting guidance for searches is diverse; Sampson et al.\ found that of eleven instruments only one item appeared in all eleven \pc{rethlefsen2021}{1}{of the eleven instruments designed to help authors report literature searches well, only one item appeared in all eleven instruments}. PRISMA 2009 had three search items (7 sources and date last searched, 8 full strategy for \emph{at least one} database, 17 numbers screened with reasons); compliance was low, partly because each item bundles several components, so authors tick the box after doing one \pc{rethlefsen2021}{2}{authors may feel if they completed one of the components of the item that they can check off that item altogether}. Librarian involvement correlates with reproducible searches, yet even librarian-authored reviews contain reproducible searches only 64\,\% of the time \pc{rethlefsen2021}{2}{reviews where librarians are authors still include reproducible searches only 64\% of the time}; little guidance existed for web sites or registries. Four aims: extensive guidance, a verification checklist for authors, editors and reviewers, interdisciplinary applicability, complement PRISMA \pc{rethlefsen2021}{2}{To develop an interdisciplinary checklist applicable to all method-driven literature searches for evidence synthesis}.

\cf{Development} Protocol registered on the EQUATOR Network, following the guidance for developers of reporting guidelines. 405 candidate items from 61 sources, consolidated to 123 \pc{rethlefsen2021}{2}{We identified 405 potential items relevant to reporting searches in systematic re- views from 61 sources}; three-round Delphi survey (round 1 rating on a 4-point scale, items retained if at least 70\,\% rated 3 or 4 and mean $\geq$ 3.25; round 2 picking the 25 most essential of 53; round 3 the same plus preferred location); 163 experts invited from 23 countries, 52 answered (32\,\%), 35 completed rounds 2 and 3 \pc{rethlefsen2021}{3}{We received 52 responses (32\% response rate) to the first survey, and of these, 35 (67\% response rate) completed both surveys two and three}; consensus conference in May 2016 with 38 attendees; 30 items remained, 3 added back; public draft 20 March 2019 and 358 comments from 22 individuals and organisations \pc{rethlefsen2021}{3}{358 separate comments were received from 22 individuals and organizations}; coordination with the PRISMA 2020 lead. All materials are on OSF.

\cf{The 16 items (Table 1)} \emph{Information sources and methods}: 1 database name with platform; 2 multi-database searching (name the platform and every database); 3 study registries; 4 online resources and browsing; 5 citation searching; 6 contacts; 7 other methods. \emph{Search strategies}: 8 full strategies, ``copied and pasted exactly as run''; 9 limits and restrictions with justification; 10 search filters (cite them); 11 prior work (cite reviews whose strategies were reused); 12 updates; 13 dates of searches. \emph{Peer review}: 14. \emph{Managing records}: 15 total records per source; 16 deduplication method and software \pc{rethlefsen2021}{4}{Include the search strategies for each database and information source, copied and pasted exactly as run}. The checklist is for reporting, not conduct, and ``not every systematic review will make use of all of the items'' \pc{rethlefsen2021}{4}{Not every systematic review will make use of all of the items in the Information Sources and Methods section of the checklist}; it may also document a planned search in a protocol (with PRISMA-P).

\cf{Key points of the elaboration} \emph{Item 1}: the same database behaves differently by platform (field codes, phrase searching, truncation, extra data), so name database \emph{and} platform \pc{rethlefsen2021}{5}{It is important that authors indicate not only the database, but the platform through which the database was searched}. \emph{Item 4 (web and grey sources)}: searching beyond databases reduces bias because databases skew to English-language positive results; for web search engines list the sites with addresses, state whether personalisation was reduced (for example a private browser window) and, if only the first results were screened, give the cap and the rationale \pc{rethlefsen2021}{8}{If authors used a general search engine, authors should declare whether steps were taken to reduce personalization bias}. The example reports a Google query, the date and ``the first 20 results were selected and screened'' \pc{rethlefsen2021}{7}{The first 20 results were selected and screened}. \emph{Item 7}: even non-replicable methods (personal files, ``related articles'' tools) should be declared, naming the tool and the seed articles. \emph{Item 8}: report all strategies, including web and registry searches, preferably in a permanent repository with a DOI (institutional repository, OSF, figshare, Zenodo) because journal supplements disappear \pc{rethlefsen2021}{10}{supplements may disappear}; a worked non-database entry records date, portal and URL, search terms and a note such as ``First 5 pages screened on title (n=50 records)'' \pc{rethlefsen2021}{10}{Notes: First 5 pages screened on title (n=50 records)}. \emph{Item 9}: database limits are ``not recommended'' in systematic reviews because they depend on indexing and risk missing recent records; any limit needs a justification, and limits must not be confused with filters or with inclusion criteria \pc{rethlefsen2021}{12}{Database limits are not recommended for use in systematic reviews, due to their fallibility}. \emph{Item 12--13}: searches should ideally end less than six months before publication; give the date of each search and describe changes on re-running (the ERIC example had to change the strategy because the database algorithm changed) \pc{rethlefsen2021}{13}{The last search in a review should be conducted ideally less than 6 months before publication}. \emph{Item 15--16}: record counts per source, because a drastically different number on repeating a search signals errors or a changed database; describe deduplication software or manual steps \pc{rethlefsen2021}{14}{If instead, the searcher locates a drastically different number of results than reported in the original review, this can be indicative of errors in the published search}.

\cf{Glossary and repositories} Table 2 distinguishes \emph{limits} (built-in restrictions) from \emph{filters} (validated combinations of terms, also called hedges) and from screening \emph{inclusion criteria}; defines platform, field code, record, search strategy, sensitivity and specificity \pc{rethlefsen2021}{6}{Limits are different from filters (see above) and are also not the inclusion/exclusion criteria used in the screening process}. Table 3 recommends depositing strategies and, optionally, files of retrieved, deduplicated and included references (abstracts removed for copyright), with a README \pc{rethlefsen2021}{11}{Authors who wish to share these files should note that abstracts are copyrighted materials and should be removed from files before sharing them publicly}.

\cf{Authors' stated limits} PRISMA-S does not address text-mining methods used to \emph{build} a search and does not require the \emph{rationale} for term choices; it is not a tool to guide conduct or to evaluate the quality of a search or review \pc{rethlefsen2021}{15}{PRISMA-S does not address using text min- ing or text analysis methods to create the search}; new information sources may require adaptations \pc{rethlefsen2021}{16}{is not intended as a tool to either guide conduct of a systematic review or to evaluate the quality of a search or a systematic review}.

\cf{Limits (card)} Biomedical, database-centred origin; the Delphi had a 32\,\% first-round response rate, mostly librarians and review methodologists, so it reflects their practice; no evidence that using the checklist improves reproducibility (the paper cites only correlations for librarian involvement); app-store search is not covered: store listings change daily, rankings are personalised and region-dependent, and a store query cannot be re-run to the same result, so the advice on dates, personalisation and result caps applies only by analogy.

\cf{Use for this thesis} (1) A checklist for the search section of the thesis and for the app-selection protocol (stage A): store and platform names, country and language setting, device or emulator, query terms, ranking order used, number of results examined, date and time of capture, filters applied (category, rating, installs), justification for every restriction, and counts before and after deduplication across stores. Items 1, 4, 8, 9, 13, 15, 16 map directly. (2) Items 5, 7 and 11 fit the literature side: citation searching in the reading corpus, personal files, and reuse of search strings from earlier app-store studies (cite them). (3) A supplementary archive in the project (search log, raw lists, README) answers item 8 and the availability item of PRISMA 2020. (4) The thesis should state plainly that no checklist was completed unless it is, and name the items adapted. (5) Pairs with PRISMA 2020 (Page et al.), PRISMA-ScR (Tricco et al.) and the scraper output folders under \texttt{scripts/}.
\end{papercard}
